\section{MoE Model Selection for MAgHARCM (FreeToken-Inspired)}
\label{app:freetoken}

\subsection{Problem}

MAgHARCM currently drives translation through two local Ollama models invoked
from \texttt{internal/llm/llm.go} (lines 20--36): a \textbf{reasoning} planner
and a \textbf{coding} translator. \texttt{config.yml:12} pins
\texttt{reasoning: "gpt-oss:20b"} --- a dense $\sim$20B model. As the project
corpus grew (Sample 2: 38 files / 5,625 LoC of \texttt{oxidizer/stats}; Sample 3:
15k LoC test data; Sample 4: 100+ Java files), the dense planner hit the VRAM
wall on consumer-class hardware: a 20B dense Q4\_K\_M runs $\sim$12--13 GB on
disk and needs $\sim$14--16 GB VRAM for the KV cache at 8k context, leaving no
headroom for the 4B coding model running in parallel.

The pattern that resolves this is the \textbf{Mixture-of-Experts (MoE)}
total-params-in-RAM + active-params-in-GPU split. Total parameters can be 8B,
35B, 290B+ because they live in host RAM and stream over PCIe; active parameters
per token are 1--3B and fit on a single consumer GPU. The same insight motivated
FreeToken (FlashML, UC Berkeley + UT Austin, August 2026).

\subsection{FreeToken pattern (what it is and isn't)}

FreeToken is an edge-native MoE serving engine --- not a new model architecture.
Verified claims from the official README
(\url{https://github.com/FlashML-org/FreeToken/blob/main/README.md}):

\begin{itemize}
\item ``Unlock datacenter-class intelligence on the hardware you already own ---
  Run \textbf{290B+ frontier MoE models locally} on your gaming PC at blistering
  interactive speeds.'' [\texttt{MEASURED} --- README banner]
\item ``\textbf{Fast Edge-Native Runtime}: Provides efficient MoE serving with
  \textbf{bandwidth-adaptive CPU--GPU co-execution ($q^\star$ policy)}, full-layer
  double-buffered prefill streaming, global LRU expert caching,
  graph-compatible execution, and the FTW fast weight format.''
  [\texttt{MEASURED} --- README ``About'' \S1]
\item ``\textbf{Broad MoE \& Ecosystem Support}: Supports frontier open-weight MoE
  models (e.g., DeepSeek-V4-Flash, Qwen3.6-35B-A3B, GLM-5.2)'' [\texttt{MEASURED}]
\end{itemize}

Concrete numbers from the FreeToken model table (README + paper
arXiv:2608.16157):

\begin{itemize}
\item \textbf{Qwen3.6-35B-A3B} (35B total, A3B = 3B active) --- runs on \textbf{8 GB VRAM
  laptop (e.g.\ RTX 4060)} at \textbf{$\sim$39.3 tok/s}.
\item \textbf{DeepSeek-V4-Flash} (284B total) --- 32 GB VRAM desktop (RTX 5090) at
  $\sim$22--25 tok/s.
\item \textbf{GLM-5.2} (\textbf{753B} total) --- \textbf{single workstation GPU}
  (e.g.\ 96 GB VRAM) at $\sim$14.9 tok/s.
\end{itemize}

What FreeToken is \textbf{not}: it is not a new routing algorithm. It is a
serving runtime. The $q^\star$ policy decides per-layer whether to compute the
active expert on GPU (fast, VRAM-bound) or CPU (slow, RAM-bound) based on
measured PCIe bandwidth. Dense models get no benefit (no experts to cache).

\subsection{Local Ollama candidates}

Numbers verified against the live local \texttt{ollama list} snapshot and
Ollama library pages as of 2026-08-30.

\begin{table}[h]
\centering
\begin{tabular}{llllll}
\toprule
Tag & Total & Active & Type & Disk (local) & Fit \\
\midrule
\texttt{lfm2.5:8b-a1b-q4\_K\_M} & \textbf{8.3B} & \textbf{$\sim$1.5B} & MoE & \textbf{5.2 GB} & Best fit \\
\texttt{qwen3.8:27b} & \textbf{27.8B} & 27.8B & dense & \textbf{17 GB} & Baseline dense alt \\
\texttt{qwen2.5-coder:7b-base-q5\_0} & 7.6B & 7.6B & dense & 5.3 GB & Small dense fallback \\
\texttt{mistral:7b-instruct-v0.2-q5\_0} & 7.3B & 7.3B & dense & 5.0 GB & Small dense fallback \\
\texttt{gpt-oss:20b} (current default) & $\sim$21B & $\sim$21B & dense (MoE variant) & $\sim$12 GB & \textbf{Current pin} \\
\bottomrule
\end{tabular}
\caption{Local Ollama model candidates for MAgHARCM reasoning slot.}
\end{table}

\texttt{lfm2.5:8b-a1b-q4\_K\_M} is the MoE option that delivers 8B-class
reasoning capacity while activating only $\sim$1.5B params/token --- the exact
total/active split FreeToken targets. The naming convention is
\texttt{<total>B-a<active>B}; the ``a1b'' suffix = 1B active (Liquid AI rounds
to 1 in the tag, the model card reports 1.5B).

\subsection{Recommendation}

Replace \texttt{reasoning: "gpt-oss:20b"} in \texttt{config.yml:12} with
\textbf{\texttt{reasoning: "lfm2.5:8b-a1b-q4\_K\_M"}}.

Rationale:
\begin{enumerate}
\item \textbf{VRAM headroom.} Active params ($\sim$1.5B Q4 $\approx$ 1.0 GB on
  the GPU) plus KV cache leaves the consumer GPU free for the 2.5 GB coding
  model and $\sim$5 GB of working memory.
\item \textbf{Speed.} FreeToken benchmarks show 35B-A3B at $\sim$39 tok/s on
  8 GB VRAM; LFM2.5-8B-A1B is the smaller sibling and runs much faster than that
  even in stock Ollama.
\item \textbf{Tool-calling strength.} Liquid AI ships LFM2.5-8B-A1B with
  day-one tool-calling and reasoning (``thinking'') support --- directly aligned
  with MAgHARCM's planner/verifier agent loop.
\item \textbf{Disk/RAM budget.} 5.2 GB on disk; $\sim$8 GB peak resident RAM.
  Total weight on RAM ladder, active slice on GPU --- the FreeToken pattern
  without needing FreeToken itself (Ollama implements the same LRU expert cache
  for small MoE models in-process).
\end{enumerate}

\textbf{Validation gate.} Measure first on Sample 1
(\texttt{assets/samples/oxidizer/stats/go}, 4 files, $\sim$280 LoC). Capture
two metrics:
\begin{itemize}
\item Inference latency (sec/sample, log to \texttt{run-2026-08-30-sN.log}).
\item Quality: pass/fail of the build/test cascade (ReCodeAgent 3-tier
  validation flow from PRIM-06).
\end{itemize}

If the MoE planner shows quality regression on Sample 1, fall back to
\textbf{\texttt{qwen2.5-coder:7b-base-q5\_0}} or
\textbf{\texttt{mistral:7b-instruct-v0.2-q5\_0}} --- same disk/RAM class as the
MoE choice, dense routing, well-understood behaviour. The \texttt{qwen3.8:27b}
17 GB dense option is held in reserve only if the small dense fallbacks fail and
the host has $\geq$24 GB RAM to spare.

\subsection{Implementation}

\subsubsection{YAML patch (exact)}

\begin{verbatim}
--- a/config.yml
+++ b/config.yml
@@ -9,7 +9,7 @@ translation:
     toolchain: "cargo"
   models:
-    reasoning: "gpt-oss:20b"
+    reasoning: "lfm2.5:8b-a1b-q4_K_M"
     coding: "hf.co/unsloth/Qwen3-4B-Instruct-2507-GGUF:UD-Q4_K_XL"
     ollama_url: "http://localhost:11434"
\end{verbatim}

No code change in \texttt{internal/llm/llm.go} --- it is a pure string lookup
(\texttt{ollama.NewChatModel} with \texttt{Model: cfg.ReasoningModel} at line 22).

\subsubsection{Expected VRAM/RAM budget}

\begin{table}[h]
\centering
\begin{tabular}{lllll}
\toprule
Slot & Model & Disk & Active VRAM & Peak RAM \\
\midrule
Reasoning & \texttt{lfm2.5:8b-a1b-q4\_K\_M} & 5.2 GB & $\sim$1.0 GB & $\sim$6 GB \\
Coding & \texttt{hf.co/unsloth/Qwen3-4B-Instruct-2507-GGUF:UD-Q4\_K\_XL} & 2.5 GB & $\sim$1.6 GB & $\sim$3 GB \\
Ollama runtime + KV (8k ctx) & --- & --- & $\sim$1.5 GB & $\sim$2 GB \\
\midrule
\textbf{Total} & & \textbf{7.7 GB} & \textbf{$\sim$4.1 GB} & \textbf{$\sim$11 GB} \\
\bottomrule
\end{tabular}
\caption{Expected VRAM/RAM budget after the reasoning-slot swap.}
\end{table}

Fits comfortably on a 12 GB consumer GPU (e.g.\ RTX 4070) with the coding
model resident. On 8 GB VRAM hardware, the coding model will spill to RAM but
the reasoning MoE stays on-GPU --- the inverse of the FreeToken 35B-on-8GB demo.

\subsubsection{Rollback plan}

\begin{enumerate}
\item Revert \texttt{config.yml:12} to \texttt{gpt-oss:20b} (or whichever prior
  pin was in effect before this change). One-line edit.
\item No persistent state changes --- Ollama keeps pulled weights in
  \texttt{\$OLLAMA\_MODELS} and the swap is non-destructive.
\item Translation runs are isolated per \texttt{config.yml} run; no
  half-migrated artifacts in \texttt{.artifacts/stats/rust/} because the
  planner prompt template does not cache model identity in the output.
\end{enumerate}

\subsection{Caveats}

\begin{itemize}
\item \textbf{MoE routing quality on small total-param models.} Liquid AI's
  8B-A1B is among the smallest MoE models shipping as a production planner.
  Routing quality on sub-8B MoE models has historically lagged dense
  equivalents on long-horizon planning tasks. The Sample 1 quality gate exists
  specifically to surface this before committing to a multi-sample sprint.
\item \textbf{Active-param count is approximate.} The \texttt{q4\_K\_M} blob on
  disk reports \textbf{8.5B} as the total weight count (rounding in GGUF
  metadata); the \texttt{a1b} tag implies $\sim$1B active per token, but the
  Liquid AI model card reports \textbf{$\sim$1.5B active}. Use 1.5B for VRAM
  planning, 1B as the lower bound for headroom estimates.
\item \textbf{Hybrid architecture.} LFM2.5-8B-A1B is a \emph{hybrid} MoE: 18
  convolution blocks + 6 GQA layers (Liquid AI docs). The non-pure-transformer
  blocks behave differently under long-context KV-cache pressure than pure
  transformer MoEs like Qwen3.6-35B-A3B. Monitor for context-length regressions
  even if short-context Sample 1 passes.
\item \textbf{FreeToken $\neq$ Ollama.} The $q^\star$ bandwidth-adaptive split
  is a FreeToken feature. Stock Ollama uses a simpler expert-offloading scheme.
  For MAgHARCM's target scale (sub-10B MoE), Ollama's scheme is adequate. If we
  later adopt a 35B+ MoE planner, \textbf{FreeToken becomes the deployment
  surface, not Ollama} --- and the Linux-x86\_64 + NVIDIA-RTX-only constraint
  becomes binding.
\item \textbf{Concurrent model residency.} Ollama's \texttt{num\_parallel} and
  \texttt{OLLAMA\_NUM\_GPU} knobs affect whether both models can stay resident.
  For the budget above to hold, set \texttt{OLLAMA\_NUM\_GPU=99} (let Ollama
  layer both models) and keep \texttt{num\_parallel=1} so KV caches don't
  fragment.
\end{itemize}
